\documentclass[12pt,a4paper,fleqn]{article}

%------------------------------------------------
% Кодировки и шрифты
%------------------------------------------------
\usepackage[utf8]{inputenc}
\usepackage[T2A]{fontenc}
\usepackage[russian]{babel}
\usepackage{amsmath,amssymb,mathtools}

%------------------------------------------------
% Оформление, отступы и выравнивание по левому краю
%------------------------------------------------
\usepackage{geometry}
\geometry{left=20mm,right=20mm,top=20mm,bottom=20mm} 

\usepackage{setspace}
\onehalfspacing 

\raggedright 
\setlength{\parindent}{0pt} 
\setlength{\parskip}{1em} 
\setlength{\mathindent}{0pt} % Формулы прижаты влево

\allowdisplaybreaks 

\begin{document}

\noindent\textbf{\Large Ответ:} Нет, не обязательно (при $\deg P \ge 5$).

\vspace{0.5cm}
\noindent\textbf{1. Случай $\deg P = 3$ (экстремумы всегда есть):}

\begin{flalign*}
& \deg P = 3 \implies \deg P' = 2 & \\[10pt]
& P'(x_1) = P'(x_2) = 0 \ (x_1 \neq x_2) \implies x_1, x_2 \text{ --- простые корни } (\text{кратности } 1) & \\[10pt]
& \implies P'(x) \text{ меняет знак при переходе через } x_1 \text{ и } x_2 \implies \text{обе точки являются экстремумами} &
\end{flalign*}

\vspace{0.5cm}
\noindent\textbf{2. Контрпример для $\deg P = 5$:}

\begin{flalign*}
& P(x) = \frac{x^5}{5} - \frac{2x^3}{3} + x \quad (\deg P = 5 \text{ --- нечетная}) & \\[10pt]
& P'(x) = x^4 - 2x^2 + 1 = (x^2 - 1)^2 = (x-1)^2(x+1)^2 & \\[10pt]
& P'(x) = 0 \implies x_1 = -1, \quad x_2 = 1 & \\[10pt]
& P'(x) \ge 0 \ \forall x \in \mathbb{R} \implies \text{знак не меняется} \implies \text{экстремумов нет} &
\end{flalign*}

\vspace{0.5cm}
\noindent\textbf{3. Проверка контрпримера через производные высших порядков:}

\begin{flalign*}
& P''(x) = 4x^3 - 4x = 4x(x^2 - 1) \implies P''(\pm 1) = 0 & \\[10pt]
& P'''(x) = 12x^2 - 4 \implies P'''(\pm 1) = 8 \neq 0 & \\[10pt]
& \implies x_1, x_2 \text{ --- точки перегиба} &
\end{flalign*}

\vspace{0.5cm}
\noindent\textbf{4. Общий случай для произвольной нечетной степени $\deg P = 2k+1 \ge 5$:}

\begin{flalign*}
& \deg P = 2k+1 \ (k \ge 2) \implies \deg P' = 2k & \\[10pt]
& \text{Пусть } P'(x) = a(x - x_1)^{m_1}(x - x_2)^{m_2}, \text{ где } m_1, m_2 \text{ --- четные, } m_1 + m_2 = 2k & \\[10pt]
& \implies P'(x) \text{ не меняет знак при переходе через } x_1, x_2 \implies \text{экстремумов нет} &
\end{flalign*}

\end{document}
